Continuous symmetries of dynamical systems are often discovered from data as
the nullspace of a linear operator. With noisy data, a symmetry is declared
whenever a singular value looks small. This has no error control and points
the wrong way: failing to reject a symmetry is not evidence for it. We recast
symmetry discovery as an equivalence test: a generator is certified as an
approximate symmetry only when the data rule out a violation above a stated
tolerance. For vector-field models linear in their coefficients, such as
polynomial or dictionary models, a relative measure of symmetry violation is a
ratio of norms of linear images of the coefficients, so one confidence region
yields bounds for every candidate generator at once. On ten systems with
exactly known symmetries and 72,000 simulated data sets, the procedure never
certified a false symmetry, whereas a fixed threshold at the same tolerance did
so in 0.5\% of runs on average and 63\% in its worst setting. The price is
fewer detections (48.5\% against 78.4\%). The guarantee requires a model class
containing the true system, exactly observed states and a confidence region
matched to the error structure; violating the first two can manufacture false
symmetries.
